\section{Reproducción y despliegue de la aplicación}
\label{anexo:reproduccion}

Este anexo describe los pasos para desplegar el sistema completo —servidor de
seguimiento MLflow, \textit{backend} FastAPI con los agentes LangGraph y la interfaz
web Next.js— en una máquina nueva. Se ofrecen dos vías: el \textbf{despliegue con
Docker} (recomendado, sin necesidad de instalar Python ni Node) y el \textbf{despliegue
local} para desarrollo. Todo el código y las dependencias fijadas (\texttt{uv.lock},
\texttt{package.json}) están versionados en el repositorio, lo que garantiza
instalaciones deterministas.

\subsection{Requisitos previos y variables de entorno}
\label{anexo:repro-requisitos}

La única dependencia obligatoria para la vía Docker es
\href{https://www.docker.com/products/docker-desktop/}{Docker~Desktop}. Para la vía
local se requiere además \textbf{Python~$\geq$~3.12} con el gestor
\href{https://docs.astral.sh/uv/}{UV} y \textbf{Node.js~$\geq$~20}.

La configuración se lee de un fichero \texttt{.env} en la raíz del proyecto. Para
crearlo, copie la plantilla incluida \texttt{.env.example} a un nuevo fichero
\texttt{.env} y adáptelo rellenando al menos la variable \texttt{OPENAI\_API\_KEY}.

La Tabla~\ref{tab:repro-env} resume las variables relevantes. La única imprescindible
es la clave de la API del proveedor de modelos; el resto tienen valores por defecto
sensatos.

\begin{table}[H]
\centering
\caption{Variables de entorno principales (\texttt{.env}).}
\label{tab:repro-env}
\small
\setlength{\tabcolsep}{5pt}
\begin{tabular}{lll}
\toprule
\textbf{Variable} & \textbf{Obligatoria} & \textbf{Descripción / valor por defecto} \\
\midrule
\texttt{OPENAI\_API\_KEY}        & Sí & Clave de la API de OpenAI (o \textit{token} de GitHub) \\
\texttt{OPENAI\_BASE\_URL}       & No & \textit{Endpoint} alternativo; p.~ej.\ GitHub~Models
                                       (Sección~\ref{anexo:repro-github}) \\
\texttt{MLFLOW\_TRACKING\_URI}   & No & \texttt{http://mlflow:5000} (Docker) /
                                       \texttt{sqlite:///./mlflow.db} (local) \\
\texttt{MLFLOW\_EXPERIMENT\_NAME}& No & \texttt{mlops-agents} \\
\texttt{LOG\_LEVEL}              & No & \texttt{INFO} \\
\bottomrule
\end{tabular}
\end{table}

\noindent
Los modelos se asignan por nodo en los ficheros de \textit{prompt} de cada agente:
\texttt{gpt-5.4-mini} para el validador de datos y el planificador (que exigen un
seguimiento preciso de instrucciones) y \texttt{gpt-5.4-nano} para el redactor de
informes. Como alternativa gratuita a la API de pago, el sistema también puede
ejecutarse contra GitHub~Models (Sección~\ref{anexo:repro-github}).

\subsection{Ejecución gratuita con GitHub Models}
\label{anexo:repro-github}

Para no obligar a contratar la API de pago de OpenAI, el sistema conserva la posibilidad
de ejecutarse contra \href{https://github.com/marketplace/models}{GitHub~Models}, el
servicio gratuito de inferencia de GitHub. Como el acceso al LLM se realiza mediante un
cliente compatible con la API de OpenAI, no se requiere ningún cambio en el código del
sistema; basta con:

\begin{enumerate}
  \item Generar un \textit{token} personal de GitHub (sin permisos especiales) y usarlo
        como valor de \texttt{OPENAI\_API\_KEY}.
  \item Apuntar el cliente al \textit{endpoint} de GitHub~Models mediante la variable de
        entorno \texttt{OPENAI\_BASE\_URL}.
  \item Sustituir el modelo de cada agente —definido en los ficheros de \textit{prompt}
        \texttt{src/mlops\_agents/prompts/} (\texttt{data\_agent}, \texttt{planner},
        \texttt{report\_writer})— por uno disponible en el catálogo de GitHub~Models
        (p.~ej.\ \texttt{gpt-4.1-mini}).
\end{enumerate}

\noindent
La contrapartida es un límite de 150~peticiones diarias por modelo en la capa gratuita y
una fiabilidad algo menor en los agentes que exigen un seguimiento preciso de
instrucciones —el validador de datos y el planificador—. Por ello, para obtener
resultados consistentes y reproducir las cifras de la Sección~\ref{sec:resultados} se
recomienda la API de OpenAI; GitHub~Models se ofrece como vía de exploración sin coste.

\subsection{Despliegue con Docker (recomendado)}
\label{anexo:repro-docker}

Desde una instalación limpia:

\begin{lstlisting}[language=bash]
git clone https://github.com/mariovuniovi/tfg_agentes_ia_langraph.git
cd tfg_agentes_ia_langraph
# copie .env.example a un fichero .env y rellene OPENAI_API_KEY
docker compose up --build     # construye las imagenes y arranca los 3 servicios
\end{lstlisting}

El comando levanta tres contenedores sobre una red privada compartida
(\texttt{mlops-net}), descritos en la Tabla~\ref{tab:repro-servicios}. Una vez en marcha,
la aplicación es accesible en \url{http://localhost:3000}.

\begin{table}[H]
\centering
\caption{Servicios desplegados por \texttt{docker compose}.}
\label{tab:repro-servicios}
\small
\setlength{\tabcolsep}{5pt}
\begin{tabular}{lll}
\toprule
\textbf{Contenedor} & \textbf{Puerto} & \textbf{Función} \\
\midrule
\texttt{mlops-mlflow}   & 5000 & Servidor de seguimiento MLflow (\textit{runs}, métricas,
                                 artefactos y registro de modelos). Persiste en el volumen
                                 \texttt{mlflow\_data}. \\
\texttt{mlops-api}      & 8000 & \textit{Backend} FastAPI + agentes LangGraph. Invoca al
                                 LLM, ejecuta el \textit{pipeline} y expone la API REST/SSE. \\
\texttt{mlops-frontend} & 3000 & Interfaz Next.js. Consume el \textit{backend} en
                                 \texttt{http://localhost:8000}. \\
\bottomrule
\end{tabular}
\end{table}

\noindent
El \textit{backend} y MLflow se comunican internamente por el nombre de red
\texttt{mlflow:5000}; el navegador usa siempre \texttt{localhost:\textit{puerto}}.
Operaciones habituales:

\begin{lstlisting}[language=bash]
docker compose up            # arrancar (sin reconstruir si no hay cambios)
docker compose up --build    # reconstruir tras cambios de codigo
docker compose down          # detener (conserva los volumenes de datos)
docker compose down -v       # detener y borrar el volumen de MLflow (reset total)
docker compose logs -f api   # ver logs de un servicio
\end{lstlisting}

\subsection{Despliegue local para desarrollo}
\label{anexo:repro-local}

Útil para iterar sin reconstruir imágenes. Se instalan las dependencias de Python con
UV y se arranca cada componente en su terminal:

\begin{lstlisting}[language=bash]
uv sync                  # instala las dependencias de Python en .venv
# copie .env.example a un fichero .env y rellene OPENAI_API_KEY

# Terminal 1 -- solo el servicio MLflow
docker compose up mlflow

# Terminal 2 -- backend FastAPI (puerto 8000)
uv run uvicorn api.main:app --reload --port 8000

# Terminal 3 -- frontend Next.js (puerto 3000)
cd frontend && npm install && npm run dev
\end{lstlisting}

\subsection{Sembrado del \textit{pool} de experiencias}
\label{anexo:repro-pool}

El planificador recupera \textit{priors} de un \textit{pool} de experiencias previas.
Para poblarlo con los \textit{datasets} de \textit{benchmark} (Anexo~\ref{anexo:datasets})
se ejecuta el \textit{runner} \textit{offline}, que descarga los conjuntos públicos
(\texttt{scikit-learn} / OpenML / Yahoo~Finance / CSV locales), entrena el ejecutor sobre
cada uno y registra un \texttt{ExperienceRecord} en
\texttt{storage/mlops\_metadata.db}:

\begin{lstlisting}[language=bash]
uv run python scripts/run_benchmark.py --trials 8
\end{lstlisting}

\noindent
Este paso es determinista y no realiza llamadas al LLM, pero requiere conexión a
internet para descargar los \textit{datasets} reales.

\subsection{Verificación}
\label{anexo:repro-tests}

La instalación puede comprobarse ejecutando la batería de pruebas, que no realiza
llamadas reales al LLM:

\begin{lstlisting}[language=bash]
uv run python -m pytest -m "not integration"   # pruebas unitarias del backend
cd frontend && npx vitest run                  # pruebas del frontend
\end{lstlisting}

\noindent
Una vez desplegado, la verificación funcional consiste en abrir
\url{http://localhost:3000}, subir un \textit{dataset} y comprobar que el \textit{pipeline}
recorre los nodos hasta la puerta HITL de despliegue; la documentación interactiva de la
API está en \url{http://localhost:8000/docs} y el registro de experimentos en
\url{http://localhost:5000}.
